\documentclass[../../main.tex]{subfiles}
\begin{document}

\begin{center}
  {\fontsize{48pt}{58pt}\selectfont \textbf{1970}\normalsize\textbf{年度}}
\end{center}

\begin{table}[htb]
  \centering
  \caption{出題テーマと難易度}
  \begin{tabular}{|c|p{6.2cm}|c|c|c|}
    \hline
    \textbf{問題番号} & \textbf{テーマ} & \textbf{思考量} & \textbf{計算量} & \textbf{総合} \\
    \hline
    第1問 & 接線の条件から求める導関数 & 易 & 易 & 易 \\
    \hline
    第2問 & 平均と階乗の不等式の帰納法 & 普 & 普 & 普 \\
    \hline
    第3問 & 空間の2直線間の対応する中点 & 易 & 易 & 易 \\
    \hline
    第4問 & 円周上の点から弦の両端までの距離の積 & 難 & 難 & 難 \\
    \hline
    第5問 & 三角関数を含む積分の極大値 & 易 & 易 & 易 \\
    \hline
    第6問 & 漸化式で定まる数列の収束 & 易 & 易 & 易 \\
    \hline
  \end{tabular}
\end{table}

\begin{table}[htb]
  \centering
  \caption{解答}
  \renewcommand{\arraystretch}{1.4}
  \begin{tabular}{|c|c|p{9cm}|}
    \hline
    \textbf{問題番号} & \textbf{小問} & \textbf{答え} \\
    \hline
    第1問 & -- & $g\prime(a)=0$ \\
    \hline
    第2問 & -- & $\left(\dfrac{n+1}{2}\right)^n>n!$（証明問題） \\
    \hline
    第3問 & -- & $M_1,M_2,M_3$は同一直線上にある（証明問題） \\
    \hline
    第4問 & -- & 劣弧$AB$の中点と，優弧$AB$上の2点．後者の位置は解答を参照 \\
    \hline
    第5問 & -- & $x=(2n-1)\pi$で極大値$\dfrac{2}{3}(2n-1)\pi$（$n=1,2,\ldots$） \\
    \hline
    第6問 & (1) & $1<x_n<a$（$n\ge2$） \\
    \cline{2-3}
    & (2) & $x_n-1\le\dfrac{a+1}{3}(x_{n-1}-1)$ \\
    \cline{2-3}
    & (3) & $\displaystyle\lim_{n\to\infty}x_n=1$ \\
    \hline
  \end{tabular}
\end{table}

\paragraph{試験形態}
6問．試験時間は要確認．

\end{document}
